\usepackage{fontspec}
\usepackage[brazil]{babel}
\usepackage{microtype}

% Formato físico obrigatório: ISO A5, 148 × 210 mm.
% Margens espelhadas consideram encadernação e leitura em páginas opostas.
\usepackage[
  paperwidth=148mm,
  paperheight=210mm,
  inner=18mm,
  outer=14mm,
  top=16mm,
  bottom=18mm,
  bindingoffset=3mm,
  headheight=13pt,
  headsep=5mm,
  footskip=9mm
]{geometry}

\usepackage{graphicx}
\usepackage{eso-pic}
\usepackage{booktabs}
\usepackage{longtable}
\usepackage{amsmath}
\usepackage[dvipsnames]{xcolor}
\usepackage{listings}
\usepackage[most]{tcolorbox}
\usepackage{fancyhdr}
\usepackage{titlesec}
\usepackage{caption}
\usepackage{natbib}
\usepackage{makeidx}
\usepackage{hyperref}

% Fontes presentes no TeX Live e adequadas a uma mancha de texto A5.
\setmainfont{TeX Gyre Pagella}
\setsansfont{TeX Gyre Heros}
\setmonofont{TeX Gyre Cursor}
\setlength{\emergencystretch}{1em}

\definecolor{ZustandBrown}{HTML}{433E38}
\definecolor{ZustandGold}{HTML}{D4A24C}
\definecolor{ZustandCream}{HTML}{F6F0E5}
\definecolor{CoverBlue}{HTML}{003A70}
\definecolor{CodeBackground}{HTML}{F5F5F3}
\definecolor{CodeGreen}{HTML}{356859}
\definecolor{CodePurple}{HTML}{704A8A}

\hypersetup{
  colorlinks=true,
  linkcolor=ZustandBrown,
  citecolor=CodeGreen,
  urlcolor=MidnightBlue,
  pdftitle={Estado sem Drama},
  pdfauthor={Caetano},
  pdfsubject={Zustand e decisões de estado para times pequenos}
}

\pagestyle{fancy}
\fancyhf{}
\fancyhead[LE]{\scriptsize\nouppercase{\leftmark}}
\fancyhead[RO]{\scriptsize\nouppercase{\rightmark}}
\fancyfoot[LE]{\sffamily\small\thepage}
\fancyfoot[RO]{\sffamily\small\thepage}
\renewcommand{\headrulewidth}{0.3pt}
\renewcommand{\footrulewidth}{0pt}

% Páginas de abertura de partes, sumário e capítulos não ilustrados.
\fancypagestyle{plain}{%
  \fancyhf{}
  \fancyfoot[LE]{\sffamily\small\thepage}
  \fancyfoot[RO]{\sffamily\small\thepage}
  \renewcommand{\headrulewidth}{0pt}
  \renewcommand{\footrulewidth}{0pt}
}

% Capas ilustradas exibem somente o fólio real, fora da área da arte.
\fancypagestyle{chaptercover}{%
  \fancyhf{}
  \fancyfoot[LE]{\sffamily\small\thepage}
  \fancyfoot[RO]{\sffamily\small\thepage}
  \renewcommand{\headrulewidth}{0pt}
  \renewcommand{\footrulewidth}{0pt}
}

\titleformat{\chapter}[display]
  {\sffamily\bfseries\color{ZustandBrown}}
  {\filleft\large\chaptername\ \thechapter}
  {0.7ex}
  {\titlerule\vspace{1ex}\fontsize{23}{26}\selectfont}
\titlespacing*{\chapter}{0pt}{-4mm}{9mm}
\titleformat{\section}
  {\sffamily\bfseries\large\color{ZustandBrown}}
  {\thesection}{0.6em}{}
\titlespacing*{\section}{0pt}{3.2ex plus 0.8ex}{1.2ex}

\captionsetup{
  font=small,
  labelfont={bf,sf,color=ZustandBrown},
  width=0.92\textwidth
}

\lstdefinelanguage{TypeScript}{
  keywords={const,let,function,return,export,import,from,type,interface,extends,
    async,await,if,else,true,false,undefined,string,number,boolean},
  sensitive=true,
  comment=[l]{//},
  morecomment=[s]{/*}{*/},
  morestring=[b]',
  morestring=[b]"
}
\lstset{
  language=TypeScript,
  basicstyle=\ttfamily\fontsize{7.2}{8.7}\selectfont,
  backgroundcolor=\color{CodeBackground},
  keywordstyle=\color{CodePurple}\bfseries,
  commentstyle=\color{CodeGreen}\itshape,
  stringstyle=\color{BrickRed},
  numbers=left,
  numberstyle=\tiny\color{gray},
  numbersep=5pt,
  frame=single,
  rulecolor=\color{gray!30},
  framesep=4pt,
  xleftmargin=2mm,
  xrightmargin=1mm,
  breaklines=true,
  breakatwhitespace=false,
  columns=fullflexible,
  showstringspaces=false,
  tabsize=2,
  captionpos=b
}

\newtcolorbox{historia}[1][]{
  breakable,colback=ZustandCream,colframe=ZustandGold!80!black,
  fonttitle=\sffamily\bfseries,title=Do dia a dia,#1
}
\newtcolorbox{decisao}[1][]{
  breakable,colback=blue!3,colframe=ZustandBrown,
  fonttitle=\sffamily\bfseries,title=Decisão de arquitetura,#1
}
\newtcolorbox{tarefa}[1][]{
  breakable,colback=green!3,colframe=CodeGreen,
  fonttitle=\sffamily\bfseries,title=Chegou uma tarefa,#1
}
\newtcolorbox{resumo}[1][]{
  colback=gray!4,colframe=ZustandBrown,
  fonttitle=\sffamily\bfseries,title=O que levamos deste capítulo,#1
}

% Área pautada para exercícios impressos, inspirada em um editor de código.
\newcommand{\espacoresposta}[1]{%
  \begin{tcolorbox}[
    enhanced,
    colback=CodeBackground,
    colframe=gray!45,
    boxrule=0.5pt,
    arc=1mm,
    left=2mm,
    right=2mm,
    top=2mm,
    bottom=1mm,
    title={\textcolor{BrickRed}{\textbullet}\enspace
      \textcolor{ZustandGold}{\textbullet}\enspace
      \textcolor{CodeGreen}{\textbullet}\hfill
      \textcolor{gray!75}{\scriptsize resposta.md}},
    coltitle=gray!75,
    colbacktitle=gray!12,
    fonttitle=\ttfamily\scriptsize,
    titlerule=0.3pt
  ]
    \foreach \n in {1,...,#1}{%
      \noindent
      \makebox[6mm][r]{\ttfamily\scriptsize\color{gray!70}\n}%
      \hspace{2mm}%
      \textcolor{gray!38}{\rule[-0.8mm]{\dimexpr\linewidth-11mm\relax}{0.3pt}}%
      \par\vspace{2.8mm}%
    }%
  \end{tcolorbox}%
}

% Aceita primeiro PDF e depois PNG. Só chamar quando a imagem existir.
\newcommand{\imagemcapitulo}[3]{%
  \begin{figure}[htbp]
    \centering
    \IfFileExists{images/#1.pdf}{%
      \includegraphics[
        width=\linewidth,
        height=0.78\textheight,
        keepaspectratio
      ]{images/#1.pdf}%
    }{%
      \includegraphics[
        width=\linewidth,
        height=0.78\textheight,
        keepaspectratio
      ]{images/#1.png}%
    }
    \caption{#2}
    \label{fig:#3}
  \end{figure}%
}

% Marcador editorial usado enquanto a arte ainda não foi adicionada.
\newcommand{\ilustracaopendente}[2]{%
  \begin{figure}[htbp]
    \centering
    \begin{tcolorbox}[
      width=\linewidth,
      colback=gray!4,
      colframe=gray!45,
      boxrule=0.5pt,
      arc=1mm
    ]
      \centering
      {\sffamily\bfseries Ilustração planejada\par}
      \vspace{1mm}
      {\small #2\par}
      \vspace{1mm}
      {\scriptsize\ttfamily images/#1.pdf ou images/#1.png\par}
    \end{tcolorbox}
  \end{figure}%
}

% Página pautada usada no acerto de capítulos para impressão frente e verso.
% Em vez de deixar o verso vazio, oferece uma área útil para anotações.
\newcommand{\paginadeanotacoes}{%
  \thispagestyle{plain}%
  \vspace*{5mm}%
  {\centering\sffamily\bfseries\large Anotações\par}%
  \vspace{3mm}%
  {\centering\small Registre dúvidas, decisões e ideias para experimentar no DevPanel.\par}%
  \vspace{8mm}%
  \begingroup
    \color{gray!45}%
    \foreach \linha in {1,...,18}{%
      \noindent\rule{\textwidth}{0.35pt}\par\vspace{6.2mm}%
    }%
  \endgroup
  \clearpage
}

% Página de abertura ilustrada. A arte fica dentro da zona segura de impressão:
% 12 mm no lado interno, 8 mm no externo, 14 mm abaixo e mais de 10 mm acima.
% A faixa inferior original é recortada para remover fólio e título incorporados.
\newcommand{\capitulocapa}[3][0]{%
  \clearpage
  \ifodd\value{page}\else\paginadeanotacoes\fi
  \refstepcounter{chapter}%
  \phantomsection
  \addcontentsline{toc}{chapter}{\protect\numberline{\thechapter}#3}%
  \chaptermark{#3}%
  \thispagestyle{chaptercover}%
  \AddToShipoutPictureBG*{%
    \AtPageLowerLeft{%
      \put(\LenToUnit{12mm},\LenToUnit{14mm}){%
        \includegraphics[
          width=128mm,
          height=182mm,
          keepaspectratio,
          trim=0 #1bp 0 0,
          clip
        ]{#2}%
      }%
    }%
  }%
  \null
  \clearpage
}

% Capa geral: fundo até o corte e arte reduzida para preservar a zona segura.
\newcommand{\livrocapa}[1]{%
  \begin{titlepage}
    \thispagestyle{empty}%
    \AddToShipoutPictureBG*{%
      \AtPageLowerLeft{\color{CoverBlue}\rule{\paperwidth}{\paperheight}}%
      \AtPageLowerLeft{%
        \put(0,\LenToUnit{4mm}){%
          \makebox[\paperwidth][c]{\includegraphics[height=202mm]{#1}}%
        }%
      }%
    }%
    \null
  \end{titlepage}%
}

% Contracapa geral com o mesmo tratamento da capa, sem fólio ou cabeçalho.
\newcommand{\livrocontracapa}[1]{%
  \clearpage
  \thispagestyle{empty}%
  \AddToShipoutPictureBG*{%
    \AtPageLowerLeft{\color{CoverBlue}\rule{\paperwidth}{\paperheight}}%
    \AtPageLowerLeft{%
      \put(0,\LenToUnit{4mm}){%
        \makebox[\paperwidth][c]{\includegraphics[height=202mm]{#1}}%
      }%
    }%
  }%
  \null
  \clearpage
}

\newcommand{\arquivo}[1]{\texttt{#1}}
\newcommand{\codigo}[1]{\texttt{#1}}
\makeindex
